\begin{table}[tbp]\centering\small
\caption{Related numerical constructions and the transfer problem considered here.
The rows compare their inputs and objectives; they are not a ranking of accuracy.}
\label{tab:related-methods}
\begin{tabular}{@{}>{\raggedright\arraybackslash}p{0.17\linewidth}
                  >{\raggedright\arraybackslash}p{0.23\linewidth}
                  >{\raggedright\arraybackslash}p{0.25\linewidth}
                  >{\raggedright\arraybackslash}p{0.25\linewidth}@{}}
\toprule
Construction & Input and output & Main property & Relation to the present problem\\
\midrule
Conservative remapping~\citep{ullrich}
& Discretised scalar fields on a source mesh to fields on a target mesh.
& Global integral and constant fields; monotonicity for suitable maps.
& Mesh-to-mesh transfer of scalar fields, rather than liquid velocity means from boundary data.\\[5pt]
Curl-Flow~\citep{curlflow}
& Discretely divergence-free staggered velocities to a continuous velocity interpolant.
& Pointwise incompressibility and solid-wall non-penetration.
& Reconstructs a vector potential from grid velocities; the target is pointwise interpolation.\\[5pt]
High-order embedded-boundary Stokes method~\citep{overton}
& Cell data and boundary conditions to a finite-volume time evolution.
& High-order spatial operators with an approximate incompressibility projection.
& Uses geometric moments and Taylor reconstruction to discretise a flow solver.\\[5pt]
Present construction
& Potential-flow boundary data to liquid-cell velocity means on receiver polygons.
& Shared-edge liquid-moment balance; cancellation of common trace error in flat strips.
& Matches the physical and adjacent Cartesian edge averages before division by strip thickness.\\
\bottomrule
\end{tabular}
\end{table}
